\documentclass[UTF8,fontset=fandol,zihao=-4]{ctexart}

% Shared preamble for paper-reading notes.
% Compile with XeLaTeX + biber for GB/T 7714-2015 references.

\usepackage[a4paper,margin=2.5cm,headsep=0.8cm,footskip=1cm]{geometry}
\usepackage{amsmath,amssymb,mathtools,bm}
\usepackage{graphicx}
\usepackage{booktabs,tabularx,array,longtable,multirow}
\usepackage{enumitem}
\usepackage{float}
\usepackage{caption}
\usepackage{fancyhdr}
\usepackage{xcolor}
\usepackage{hyperref}
\usepackage{bookmark}
\usepackage{csquotes}
\usepackage[backend=biber,style=gb7714-2015,sorting=none]{biblatex}

\hypersetup{
  colorlinks=true,
  linkcolor=black,
  citecolor=blue!50!black,
  urlcolor=blue!60!black,
  pdfcreator={XeLaTeX},
  pdfsubject={Paper Reading Notes}
}

\ctexset{
  section = {
    name = {, .},
    number = \arabic{section},
    format = \Large\bfseries
  },
  subsection = {
    format = \large\bfseries
  },
  subsubsection = {
    format = \normalsize\bfseries
  }
}

\setlength{\parindent}{2em}
\setlength{\parskip}{0.35em}
\linespread{1.2}
\setlist[itemize]{leftmargin=2em,itemsep=0.3em,topsep=0.3em}
\setlist[enumerate]{leftmargin=2em,itemsep=0.3em,topsep=0.3em}
\captionsetup{font=small,labelfont=bf}

% Keep page numbers stable across all pages, including the title page.
\setlength{\headheight}{14pt}
\pagestyle{fancy}
\fancyhf{}
\fancyfoot[C]{\thepage}
\renewcommand{\headrulewidth}{0pt}
\renewcommand{\footrulewidth}{0pt}
\fancypagestyle{plain}{
  \fancyhf{}
  \fancyfoot[C]{\thepage}
  \renewcommand{\headrulewidth}{0pt}
  \renewcommand{\footrulewidth}{0pt}
}

\newcolumntype{Y}{>{\raggedright\arraybackslash}X}
\newcommand{\NoteField}[2]{\textbf{#1} & #2 \\}

\addbibresource{../bib/references.bib}

\title{论文阅读总索引\\\large Reading Index}
\author{Notes Subsystem}
\date{最近维护日期：\today}

\newcommand{\paperfilelink}[1]{\href{run:#1}{\nolinkurl{#1}}}
\newcommand{\paperentry}[6]{%
  \noindent
  \begin{tabularx}{\textwidth}{@{}lY@{}}
  \NoteField{标题}{#1}
  \NoteField{作者}{#2}
  \NoteField{阅读时间}{#3}
  \NoteField{一句话摘要}{#4}
  \NoteField{标签}{#5}
  \NoteField{笔记文件}{\paperfilelink{#6}}
  \end{tabularx}
  \vspace{0.4em}
  \hrule
  \vspace{0.9em}
}

\begin{document}

\maketitle

\section{使用方式}

这个索引文件用于快速浏览和维护“已读论文”的核心信息，不依赖自动扫描。后续新增论文时，直接复制一条 \verb|\paperentry{...}{...}{...}{...}{...}{...}| 并手工填写即可。

\begin{enumerate}
  \item 在 \verb|papers/| 中创建或复制笔记文件。
  \item 在 \verb|reading_index.tex| 中追加一条 \verb|\paperentry| 记录。
  \item 如果希望标题带引文，直接在第一个参数里手工写入 \verb|标题~\cite{bibkey}|。
\end{enumerate}

\section{已收录条目}

\paperentry
  {DeepSeekMath: Pushing the Limits of Mathematical Reasoning in Open Language Models~\cite{deepseekai2024math}}
  {DeepSeek-AI}
  {2026-03-22}
  {通过大规模数学语料继续预训练、专门化数据筛选和 GRPO 强化学习，DeepSeekMath 显著提升了开源 7B 模型的数学推理能力。}
  {math reasoning; continued pretraining; GRPO; open LLM; DeepSeek}
  {../papers/sample_deepseekmath_note.tex}

\paperentry
  {DeepSeek-R1: Incentivizing Reasoning Capability in LLMs via Reinforcement Learning~\cite{deepseekai2025r1}}
  {DeepSeek-AI}
  {2026-03-19}
  {通过大规模强化学习激发推理行为，并结合冷启动数据与多阶段训练，使开放模型在数学、代码和通用推理任务上逼近顶级闭源系统。}
  {reasoning LLM; reinforcement learning; GRPO; distillation; DeepSeek}
  {../papers/sample_note.tex}

\section{References}

\printbibliography[heading=none]

\end{document}
